\documentclass{amsart}
\usepackage{amsmath,amssymb,amsthm,hyperref}

\theoremstyle{definition}
\newtheorem{definition}{Definition}
\newtheorem{thesis}{Thesis}
\theoremstyle{plain}
\newtheorem{lemma}{Lemma}
\newtheorem{proposition}{Proposition}
\newtheorem{theorem}{Theorem}
\newtheorem{corollary}{Corollary}

\begin{document}

\title{Against Instrumentalism: The Irreducible Theoretical Function of
Scientific Theories}
\maketitle

\begin{abstract}
We give a structured argument settling the question posed: whether
\emph{instrumentalism}---the view that the meaning and value of a scientific
theory are \emph{exhausted} by its predictive utility as an instrument---is the
correct primary interpretation of scientific theories. We argue that it is
mistaken. The argument proceeds by first fixing precise definitions, then
establishing two asymmetries between the use of statements as instruments
(applied science) and as conjectures (pure science): an \emph{asymmetry of
truth-bearing} and an \emph{asymmetry of falsifiability}. From these we derive
that the theoretical (descriptive, truth-apt) function of a theory is not
reducible to its instrumental function, which refutes instrumentalism as a
\emph{primary} interpretation. Finally we determine, on the resulting view,
what the theoretical (non-practical) function of predictions themselves
consists in: not the practical anticipation of events, but the production of
\emph{potential falsifiers}, i.e.\ the test-statements by which a theory is
exposed to refutation and thereby treated as a description of reality.
\end{abstract}

\section{Formalization}

We first eliminate ambiguity by stating exactly what is being compared and
what counts as a successful resolution.

\begin{definition}[Statements and their two uses]\label{def:uses}
Let a \emph{statement} be any declarative sentence of a scientific theory
(a universal law such as ``all planets move on ellipses,'' a formula such as
an equation of motion, or a singular prediction such as ``the eclipse begins
at $t_0$''). A statement may be \emph{used} in two ways.
\begin{enumerate}
\item[(A)] \textbf{Instrumental use} (applied science). The statement is used
as a \emph{tool} or rule of computation: it is fed inputs (initial conditions,
measured quantities) and yields outputs (predictions) \emph{for a given
practical purpose}. On this use the statement is appraised only as
\emph{efficient} or \emph{inefficient}, \emph{applicable} or
\emph{inapplicable}, \emph{convenient} or \emph{inconvenient} to that purpose.
\item[(C)] \textbf{Conjectural use} (pure science). The statement is used as a
\emph{tentative description of, or conjecture about, reality}: an assertion put
forward as a candidate for being \emph{true} or \emph{false} of the world,
independently of any practical end.
\end{enumerate}
\end{definition}

\begin{definition}[The two positions]\label{def:positions}
\emph{Instrumentalism} is the thesis that use (A) is \emph{primary} and
exhaustive: the entire meaning and value of any scientific statement is given
by its role as an instrument for prediction, so that whatever appears to be a
function of type (C) either (i) reduces without remainder to a function of type
(A), or (ii) is a confusion to be eliminated.

The \emph{theoretical (realist) view} we shall defend is the negation of this:
scientific statements have a function of type (C)---a genuinely theoretical,
truth-apt, descriptive function---that is \emph{irreducible} to any function of
type (A).
\end{definition}

\begin{definition}[Reducibility]\label{def:reduce}
A function $F$ of statements \emph{reduces} to a function $G$ if every property
that $F$ confers on a statement, and every appraisal proper to $F$, can be
defined in, or paraphrased without loss into, the vocabulary proper to $G$. If
some property or appraisal essential to $F$ cannot be so paraphrased, $F$ is
\emph{irreducible} to $G$.
\end{definition}

The question to be settled is therefore the truth value of the following.

\begin{thesis}[To be refuted]\label{th:inst}
The conjectural function \textup{(C)} reduces to the instrumental function
\textup{(A)}; equivalently, instrumentalism gives the correct primary
interpretation of scientific theories.
\end{thesis}

\section{Two structural asymmetries}

We isolate two features that the conjectural use (C) possesses and the
instrumental use (A) lacks. By Definition~\ref{def:reduce}, exhibiting even one
appraisal essential to (C) that cannot be paraphrased in the vocabulary of (A)
suffices to refute Thesis~\ref{th:inst}. We establish two, for robustness.

\subsection{Asymmetry of truth-bearing}

\begin{lemma}[Instruments are not true or false]\label{lem:instr-notrue}
Under the instrumental use \textup{(A)}, the appraisals proper to a statement
are exhausted by the relational, purpose-indexed predicates \emph{applicable /
inapplicable}, \emph{efficient / inefficient}, \emph{convenient /
inconvenient} (relative to a fixed practical purpose). The unrelativized
predicates \emph{true} and \emph{false} are not among them.
\end{lemma}

\begin{proof}
This follows from the meaning of ``instrument'' fixed in
Definition~\ref{def:uses}(A). An instrument is appraised solely by how well it
serves an end; the appraisal is therefore ineliminably indexed to that end. A
hammer is good \emph{for} driving nails, useless \emph{for} cutting; it is
never \emph{true} or \emph{false}. The same holds for a computational rule used
as a tool: a rule of thumb, an approximation formula, or a deliberately false
idealization (e.g.\ ``treat the gas molecules as point masses,'' ``treat the
Earth as flat over short distances'') can be \emph{maximally efficient and
applicable for its purpose while being literally false}. Hence in the
vocabulary proper to (A) the question ``Is the statement true?'' is not merely
hard but \emph{undefined}: the only well-formed questions are ``True/efficient
\emph{for what}?'' Truth, being a non-relational appraisal, is not expressible
by any purpose-indexed predicate of (A). \emph{For if it were, it would inherit
the purpose-index, contradicting its non-relational character.}
\end{proof}

\begin{lemma}[Conjectures are essentially true or false]\label{lem:conj-true}
Under the conjectural use \textup{(C)}, the predicates \emph{true} and
\emph{false} are essential appraisals: to use a statement as a conjecture about
reality \emph{just is} to put it forward as something that is the case or not
the case, and so as a bearer of truth value.
\end{lemma}

\begin{proof}
Immediate from Definition~\ref{def:uses}(C): a ``tentative description of
reality'' is an assertion claiming that reality is thus-and-so. Such a claim is
constitutively the kind of thing that can match or fail to match reality, i.e.\
be true or false. Strip away truth-aptness and nothing is left of the notion of
a description \emph{of reality}; what remains is at most a rule, which is use
(A), not use (C). Thus truth-aptness is not an optional adornment of (C) but
its defining mark.
\end{proof}

\begin{proposition}[First irreducibility]\label{prop:irred1}
The conjectural function \textup{(C)} does not reduce to the instrumental
function \textup{(A)}.
\end{proposition}

\begin{proof}
By Lemma~\ref{lem:conj-true}, the appraisal ``true / false'' is essential to
(C). By Lemma~\ref{lem:instr-notrue}, this appraisal is not expressible in, and
cannot be paraphrased into, the vocabulary proper to (A), because (A) admits
only purpose-indexed appraisals while truth is non-relational. Hence an
appraisal essential to (C) cannot be paraphrased without loss into (A). By
Definition~\ref{def:reduce}, (C) is irreducible to (A).
\end{proof}

\subsection{Asymmetry of falsifiability}

The first asymmetry already settles the question, but instrumentalism's
characteristic rejoinder is that ``truth'' is metaphysical surplus and that
the only \emph{operationally} respectable virtue of a theory is its
\emph{testability}. We meet this rejoinder on its own ground by showing that
even testability behaves differently in the two uses, and in a way that again
breaks reduction.

\begin{lemma}[Instruments cannot break down in the sense of being
refuted]\label{lem:instr-norefute}
Under use \textup{(A)}, a statement can \emph{at most} reveal the limits of its
\emph{range of application}; it can never be \emph{refuted}. A clash between a
tool's output and an observation shows only that the tool has been applied
\emph{outside the range for which it is convenient}, not that the tool is
false; one simply restricts its domain of use and retains it.
\end{lemma}

\begin{proof}
By Lemma~\ref{lem:instr-notrue} an instrument bears no truth value, so it
cannot acquire the truth value \emph{false}; hence it cannot be
\emph{falsified} in the strict sense (shown to be false). When predicted and
observed results diverge, the instrumentalist appraisal is purpose-indexed
(Definition~\ref{def:uses}(A)): the verdict can only be ``inapplicable /
inefficient \emph{here},'' which is consistent with ``applicable / efficient
\emph{there}.'' Thus a divergence prompts a \emph{restriction of range}
(``Newtonian mechanics is a fine instrument at low velocities''), never an
unrestricted rejection. An instrument is therefore in principle immune to
refutation: it can be set aside as inconvenient but not convicted of error.
\end{proof}

\begin{lemma}[Conjectures can be refuted]\label{lem:conj-refute}
Under use \textup{(C)}, a statement asserted of reality \emph{can} clash with
reality and thereby be \emph{refuted}: a single observed counter-instance,
incompatible with what the conjecture asserts, establishes that the conjecture
is false.
\end{lemma}

\begin{proof}
By Lemma~\ref{lem:conj-true} a conjecture is true or false according as reality
is or is not as it asserts. If an observation reports that reality is otherwise
than the conjecture asserts, the conjecture's truth value is thereby fixed as
\emph{false}; this is refutation. Crucially, the verdict is \emph{not}
purpose-indexed and so cannot be deflected into a mere restriction of range
without \emph{abandoning the conjectural use itself}: to respond ``then the law
is merely inapplicable here'' is precisely to stop treating the statement as a
description of reality and to start treating it as a tool, i.e.\ to switch from
(C) to (A).
\end{proof}

\begin{proposition}[Second irreducibility]\label{prop:irred2}
The conjectural function \textup{(C)} does not reduce to the instrumental
function \textup{(A)}, on the ground of refutability alone.
\end{proposition}

\begin{proof}
Refutability---the capacity to be convicted of \emph{falsity} by a clash with
reality---is, by Lemma~\ref{lem:conj-refute}, available to (C). By
Lemma~\ref{lem:instr-norefute} it is unavailable to (A), which admits only
restriction of range. Refutability is an appraisal essential to the scientific
use of conjectures (an untestable conjecture is not scientifically serious),
yet it cannot be paraphrased into (A): the (A)-paraphrase of ``refuted'' is
``found inconvenient here,'' which is logically weaker (compatible with
retention) and so loses content. Hence, again by
Definition~\ref{def:reduce}, (C) is irreducible to (A).
\end{proof}

\section{The main theorem}

\begin{theorem}[Instrumentalism is mistaken as a primary
interpretation]\label{thm:main}
Thesis~\textup{\ref{th:inst}} is false. Scientific theories possess a
theoretical (descriptive, truth-apt, refutable) function that is irreducible to
their instrumental function. Therefore instrumentalism does not give the
correct primary interpretation of scientific theories.
\end{theorem}

\begin{proof}
Thesis~\ref{th:inst} asserts that (C) reduces to (A). By
Proposition~\ref{prop:irred1} (and independently by
Proposition~\ref{prop:irred2}) this is false: there are appraisals essential to
(C)---truth-bearing and refutability---that cannot be paraphrased without loss
into the vocabulary of (A). Hence (C) is a genuine, irreducible function of
scientific statements.

It remains to see that this irreducible function is \emph{primary}, i.e.\ that
instrumentalism cannot be salvaged by demoting (C) to a derivative or
dispensable role. Two observations establish this.

\emph{(i) Asymmetric dependence.} The instrumental use presupposes the
conjectural use, but not conversely. To use a formula as a reliable
\emph{tool}---to trust its outputs as predictions of what will be
observed---one must already take it (or the theory behind it) to \emph{describe
how the relevant part of the world behaves}; the tool is reliable exactly
insofar as the underlying description is approximately true. Thus the
practical, instrumental success of applied science is \emph{parasitic} on the
descriptive content supplied by pure science: applied science applies the
conjectures of pure science. The converse dependence fails---a conjecture can
be entertained, tested, and even refuted with no practical application in view
(Lemma~\ref{lem:conj-refute}; cf.\ purely theoretical predictions confirmed or
refuted long before any use). Therefore (C) is explanatorily and historically
prior to (A); a function on which the other depends, and which is irreducible
to it, cannot be the \emph{secondary} one. So if either use is ``primary,'' it
is (C).

\emph{(ii) Loss of distinctions.} An interpretation is correct only if it
preserves the appraisals that scientific practice actually makes. Pure science
distinguishes \emph{true theory from false theory} and \emph{refuted theory
from surviving theory}; by Lemmas~\ref{lem:instr-notrue} and
\ref{lem:instr-norefute} the purely instrumental vocabulary collapses both
distinctions (everything becomes ``more or less convenient, within some
range''). An interpretation that erases the very distinctions science draws
between its successes and its errors cannot be its \emph{correct primary}
interpretation. Hence instrumentalism, while a correct account of the
\emph{applied}, derivative use (A), is \emph{mistaken} as the primary
interpretation of scientific theories.
\end{proof}

\section{The theoretical function of predictions}

The problem asks, in addition, what the theoretical (non-practical) function of
\emph{predictions} consists in on the view just established. The instrumentalist
takes prediction to be the whole point of theory; on the realist view
prediction is retained but \emph{reinterpreted}, acquiring a function it does
not have on the instrumental reading.

\begin{definition}[Potential falsifier]\label{def:falsifier}
A \emph{potential falsifier} of a theory $T$ is a singular observation-statement
$f$ such that $f$ is logically incompatible with $T$ (so that, were $f$
observed, $T$ would be refuted in the sense of
Lemma~\ref{lem:conj-refute}). The class of potential falsifiers of $T$ is the
class of conceivable observations $T$ \emph{forbids}.
\end{definition}

\begin{proposition}[Predictions as potential
falsifiers]\label{prop:pred-falsifier}
A prediction derived from a theory $T$ together with initial conditions is, on
the theoretical view, an instruction for generating a potential falsifier of
$T$: it specifies in advance an observable outcome whose \emph{negation}, if
observed, would refute $T$. Its theoretical (non-practical) function is to
\emph{expose $T$ to refutation}, i.e.\ to convert $T$'s descriptive claim about
reality into a definite confrontation with reality.
\end{proposition}

\begin{proof}
Let $T$ be used conjecturally (use (C)) and let $c$ be a statement of initial
conditions. Deduce a singular prediction $p$: ``the observable outcome will be
$O$.'' Then the contrary singular statement $\bar p$: ``the observable outcome
is not $O$'' is logically incompatible with the conjunction $T\wedge c$. Given
$c$, observing $\bar p$ refutes $T$ by Lemma~\ref{lem:conj-refute}. Hence $\bar
p$ is a potential falsifier of $T$ in the sense of
Definition~\ref{def:falsifier}, and issuing the prediction $p$ is exactly the
act of singling out, in advance, an observation that $T$ forbids. The point of
issuing $p$ in pure science is therefore not to \emph{use} the outcome for any
external purpose but to \emph{stake $T$'s descriptive claim on it}: to make
$T$'s assertion about reality testable at a definite point. This is a function
of type (C), not type (A): it is not measured by convenience to a purpose but
by what it would \emph{teach us about the truth or falsity of $T$}.
\end{proof}

\begin{corollary}[The non-practical point of prediction]\label{cor:why-predict}
On the resulting view the theoretical function of a prediction is twofold and
purely epistemic:
\begin{enumerate}
\item \textbf{As a test.} A prediction is the cutting edge by which a theory is
brought into contact with reality; the more precise and the more
\emph{improbable in advance} the predicted outcome, the larger the class of
potential falsifiers it activates, and hence the more severely the theory is
tested (Proposition~\ref{prop:pred-falsifier}).
\item \textbf{As a bearer of explanatory content.} Because a refutable
prediction forbids something, it has empirical \emph{content} (it excludes
possible states of the world). A theory's descriptive, explanatory grasp of the
world is measured by the totality of what it thus forbids. Predictions are the
concrete tokens of that content.
\end{enumerate}
In neither role is the prediction valued for practical use; in both it serves
the theoretical end of learning whether, and how far, the theory's description
of reality is true.
\end{corollary}

\begin{proof}
Both clauses are immediate consequences of
Proposition~\ref{prop:pred-falsifier} together with
Definition~\ref{def:falsifier}. For~(1): the class of potential falsifiers
activated by $p$ is the class of observations incompatible with $p$ given $c$;
the more specific (the less antecedently probable) the predicted $O$, the
larger this class, hence the greater the risk to which $T$ is exposed---which is
the measure of test severity. For~(2): a statement has empirical content
exactly to the extent that it has potential falsifiers (it rules possibilities
out); by Proposition~\ref{prop:pred-falsifier} a prediction has nonempty such a
class, so it carries content, and the aggregate of a theory's forbidden
observations is its empirical content, i.e.\ the precise sense in which it
\emph{says something about the world}. Neither measure refers to any practical
purpose, so both functions are of type (C).
\end{proof}

\section{Conclusion}

\begin{theorem}[Resolution]\label{thm:resolution}
Instrumentalism is correct only as an account of the \emph{applied},
derivative use of scientific statements as tools, and is \emph{mistaken} as the
primary interpretation of scientific theories: theories possess an irreducible
theoretical (descriptive, truth-apt, refutable) function on which their
instrumental use itself depends. Correspondingly, the theoretical
(non-practical) function of predictions is not the anticipation of events for
use, but the generation of potential falsifiers---the points at which a theory
is staked against reality and thereby treated, and tested, as a description of
the world.
\end{theorem}

\begin{proof}
The first sentence is Theorem~\ref{thm:main} together with its part~(i)
(asymmetric dependence) and part~(ii) (loss of distinctions). The second
sentence is Proposition~\ref{prop:pred-falsifier} with
Corollary~\ref{cor:why-predict}. Combining them yields the stated resolution.
\end{proof}

\end{document}
